%% kapa: exam=1 round=33 year=2022 period=1 form=A subject=real_estate subject_ko=부동산학원론 seq=28 number=108
%% source: Q-Net 감정평가사 기출문제 (cst00309) artlSeq=5209997 file=raw/1008/5209997/2022년도 제33회 감정평가사 1차시험 1교시(민법,경제학,부동산학).pdf
%% conversion: kordoc 4.24.51 -> tools/gen1.py (auto); review=proofread-sonnet; flags=none
\begin{question}{108}
에스크로우(Escrow)에 관한 설명으로 옳지 않은 것은?
\begin{choices}
  \item 부동산매매 및 교환 등에 적용된다.
  \item 권리관계조사, 물건확인 등의 업무를 포함한다.
  \item 매수자, 매도자, 저당대출기관 등의 권익을 보호한다.
  \item 은행이나 신탁회사는 해당 업무를 취급할 수 없다.
  \item 에스크로우 업체는 계약조건이 이행될 때까지 금전·문서·권원증서 등을 점유한다.
\end{choices}
\end{question}
\answer{4}
